\begin{lemma}[Localization of the maximizers of a small linear perturbation]
Let $E$ be a finite-dimensional real inner-product space, let $\phi:E\to\R$, let
$z_0\in E$ and let $0<\rho\le r_1$. Assume that $\phi$ is continuous on
$\overline B(z_0,r_1)$ and that $z_0$ is a strict maximum point of $\phi$ there:
\[
  \phi(z)<\phi(z_0)\qquad\bigl(z\in\overline B(z_0,r_1),\ z\ne z_0\bigr).
\]
Then there exists $\delta_0>0$ with the following property. Whenever $p\in E$
satisfies $|p|\le\delta_0$ and $z\in\overline B(z_0,r_1)$ satisfies
\[
  \phi(w)+\langle p,w\rangle\le\phi(z)+\langle p,z\rangle
  \qquad\bigl(w\in\overline B(z_0,r_1)\bigr),
\]
one has $|z-z_0|<\rho$.
\end{lemma}
\begin{proof}
Note first that $r_1\ge\rho>0$, so $r_1>0$. Put
\[
  S=\overline B(z_0,r_1)\setminus B(z_0,\rho)
   =\bigl\{z\in E:\ \rho\le|z-z_0|\le r_1\bigr\}.
\]
The second description is the first one unwound: $z\in\overline B(z_0,r_1)$ says
$|z-z_0|\le r_1$, and $z\notin B(z_0,\rho)$ says $\neg(|z-z_0|<\rho)$, i.e.
$\rho\le|z-z_0|$.

\emph{Case 1: $S=\emptyset$.}
Take $\delta_0=1>0$. Let $p$ and $z$ be as in the statement; we must show
$|z-z_0|<\rho$. If not, then $\rho\le|z-z_0|$, and also $|z-z_0|\le r_1$
because $z\in\overline B(z_0,r_1)$; by the description of $S$ this puts $z\in
S$, contradicting $S=\emptyset$. (This is the only case that can occur when
$E=\{0\}$, and the hypothesis of the lemma is then vacuous; no nondegeneracy
assumption on $E$ is needed anywhere.)

\emph{Case 2: $S\ne\emptyset$.}
The set $\overline B(z_0,r_1)$ is closed and bounded in a finite-dimensional
space, hence compact (preamble), and $B(z_0,\rho)$ is open; so
$S=\overline B(z_0,r_1)\setminus B(z_0,\rho)$ is compact, being the
intersection of a compact set with a closed set. Also $S\subseteq\overline
B(z_0,r_1)$, on which $\phi$ is continuous, so $\phi$ is continuous on $S$.
A continuous real function on a nonempty compact set attains its maximum
(preamble): there is $z_m\in S$ with
\[
  \phi(z)\le m:=\phi(z_m)\qquad(z\in S).
\]

We claim $m<\phi(z_0)$. Indeed $z_m\in S\subseteq\overline B(z_0,r_1)$, and
$z_m\ne z_0$ because $|z_m-z_0|\ge\rho>0$ while $|z_0-z_0|=0$; so the strict
maximum hypothesis gives $\phi(z_m)<\phi(z_0)$, i.e. $m<\phi(z_0)$.

Put
\[
  \delta_0=\frac{\phi(z_0)-m}{2r_1},
\]
which is strictly positive because $\phi(z_0)-m>0$ and $r_1>0$.

Now let $|p|\le\delta_0$ and let $z\in\overline B(z_0,r_1)$ satisfy the
maximality hypothesis, and suppose for contradiction that $|z-z_0|\ge\rho$.
Since also $|z-z_0|\le r_1$, we get $z\in S$ and therefore
\begin{equation}\label{eq:smpl-m}
  \phi(z)\le m .
\end{equation}
The point $z_0$ lies in $\overline B(z_0,r_1)$ because $r_1\ge0$, so the
maximality hypothesis at $w=z_0$ gives
$\phi(z_0)+\langle p,z_0\rangle\le\phi(z)+\langle p,z\rangle$, that is,
\[
  \phi(z_0)-\phi(z)\le\langle p,z\rangle-\langle p,z_0\rangle
   =\langle p,z-z_0\rangle .
\]
By the Cauchy--Schwarz inequality and the hypotheses $|p|\le\delta_0$ and
$|z-z_0|\le r_1$,
\[
  \langle p,z-z_0\rangle\le|p|\,|z-z_0|\le\delta_0 r_1
   =\frac{\phi(z_0)-m}{2}.
\]
Combining this with \eqref{eq:smpl-m},
\[
  \phi(z_0)-m\le\phi(z_0)-\phi(z)\le\frac{\phi(z_0)-m}{2},
\]
whence $\phi(z_0)-m\le0$, contradicting $m<\phi(z_0)$. Therefore
$|z-z_0|<\rho$, as required.
\end{proof}
